% model_specification.tex -- Slide-to-Target: model specification (reference note, v1, 31 Oct 2026).
%
% A reference note for the team, written alongside the code (physics.py, observe.py, belief.py, decide.py,
% drake_config.py, e0_*.py).  It replaces the long-form proposal's paragraphs "Assumptions on the object",
% "Contact model" and "Fusing the prior with practice slides".  Under the 6.4210 AI policy the team writes the
% prose of the submitted proposal and report; this note supplies the equations, tables and figures around it.
% Build: pdflatex model_specification.tex (twice).  error_budget_rows.tex is generated by `python3 observe.py`.
\documentclass[10pt]{article}
\usepackage[utf8]{inputenc}
\usepackage[T1]{fontenc}
\usepackage[margin=2.2cm]{geometry}
\usepackage{amsmath,amssymb}
\usepackage{array,booktabs,longtable}
\usepackage{hyperref}
\usepackage{xcolor}
\setlength{\parskip}{4pt}
\setlength{\parindent}{0pt}
\newcommand{\dstar}{d^{\star}}
\newcommand{\code}[1]{\texttt{#1}}

\title{Slide-to-Target: model specification\\[2pt]\large reference note v1}
\author{Slide-to-Target team, MIT 6.4210}
\date{31 October 2026; Section~\ref{sec:status-nov} added 17 November}

\begin{document}
\maketitle

\noindent\fbox{\parbox{0.97\textwidth}{\small
\textbf{What this is.} The full generative model the instructor asked for: how a commanded push becomes a release
state, how a release state becomes a landing, what is measured, and how the belief over friction is updated and
used.  Sections~\ref{sec:state}--\ref{sec:ladder} are the specification; Section~\ref{sec:impl} maps every equation
to a function and a test in the repository.  Numbers marked \emph{provisional} wait on the E0 runs named beside them.
\textbf{Not submitted text.} The team writes the proposal and report prose; this note is the reference the code is
checked against.}}

\section{Overview}
\label{sec:overview}
The pipeline, one practice slide per loop:
\begin{center}\small
\begin{tabular}{@{}l@{\;$\rightarrow$\;}l@{\;$\rightarrow$\;}l@{\;$\rightarrow$\;}l@{}}
commanded push $u$ & push map (release model) & free slide & measurements: frames, rest pose, FK release point\\[2pt]
\multicolumn{4}{@{}l}{$\downarrow$ physics likelihood (model C) \quad $\leftarrow$ LLM observation model, prior}\\[2pt]
\multicolumn{4}{@{}l}{grid posterior over $\theta=\ln\mu$ \;$\rightarrow$\; decision (minimum expected loss) \;$\rightarrow$\; next push}
\end{tabular}
\end{center}
The first shot uses the LLM and the prior alone; each practice slide closes the loop once.  Everything below
holds for the main study; the fidelity ladder of Section~\ref{sec:ladder} says which assumptions are relaxed
where.

\section{State and parameters}
\label{sec:state}
\begin{itemize}
\item \textbf{Unknown, shared across a pair's slides:} $\theta=\ln\bar\mu$, the mean object--table kinetic friction,
      and strip offsets $\eta_s$ (level L2 and up).
\item \textbf{Known and fixed in the main study:} mass $m$, geometry, centre of mass at the base centroid (offset $e=0$),
      paddle--object friction $\mu_p$.
\item \textbf{Free at level L3:} $e$ and $\mu_p$.
\item \textbf{Per-slide nuisance:} release position $x_r$ and release speed $v_r$.  They are estimated from data, never
      taken from the commanded trajectory.
\end{itemize}
Under A1--A3 the stopping distance does not depend on the centre of mass: pure translation puts the friction
resultant under the COM, and a line contact can place its force centroid anywhere along the contact segment.  So
the COM is treated as known for boxes, and E0-3 tests where that fails.

\section{Assumptions, each with its check}
\label{sec:assumptions}
\begin{center}\small
\begin{tabular}{@{}llp{3.2cm}@{}}
\toprule
ID & Assumption & Checked by\\
\midrule
A1 & Flat rigid base, uniform pressure & E0-1\\
A2 & Line contact; the COM projection lies inside the paddle contact segment ($|e|<w/2$) & E0-3\\
A3 & Push along the face normal; separation at the first time the paddle decelerates faster than $\mu g$ & E0-4\\
A4 & Friction strips run across the push direction; $\mu$ does not depend on speed over 0.3--1.2 m/s & E0-6; rig $\mu(v)$\\
A5 & No tipping during the push, $(a+\mu g)h_p\le gb+a\,h_c$, or during the slide, $\mu h_c<b$ & analytic, per object\\
A6 & Drake: 1 ms step, Lagged approximation (tentative), one compliant geometry per contact pair, margin $10^{-4}$ m & E0-1, E0-2\\
\bottomrule
\end{tabular}
\end{center}
In A5, $b$ is the half-length along the push, $h_c$ the COM height and $h_p$ the paddle contact height
(\code{physics.no\_tip}).

\section{Contact specification}
\label{sec:contact}
\begin{center}\small
\begin{tabular}{@{}p{3.1cm}p{6.2cm}p{6.2cm}@{}}
\toprule
Property & Paddle--object & Object--table\\
\midrule
Geometry & Flat paddle face on the box's end face & Box base on the table plane\\
Contact type in the plane & Segment of width $w$ (line).  Discs in E0-3 give point contact on purpose & Full base patch\\
Dynamic friction & $\mu_p$ chosen (nominal 0.3), swept 0.1--0.6 in E0-3 & Measured on the rig (primary set), or Yu et al.\ (simulator-only set)\\
Static friction & Not used by Drake's discrete solver & Not used\\
Compliance & Object compliant, paddle rigid & Object compliant, table rigid\\
Hydroelastic modulus & $10^{7}$ Pa nominal (Drake default); swept $10^{5}$--$10^{7}$ & Same\\
Dissipation & Hunt--Crossley $d$ under Lagged/Similar (swept 0, 1, 10, 50 s/m); \code{relaxation\_time} only under SAP & Same\\
Margin & $10^{-4}$ m & $10^{-4}$ m\\
Contact height & $h_p$ chosen to satisfy A5 & Not applicable\\
How it affects the landing & Through release speed and yaw (E0-3, E0-4) & Directly (E0-1)\\
\bottomrule
\end{tabular}
\end{center}
Two Drake facts shape this table.  With two rigid hydroelastic geometries, the default
\code{hydroelastic\_with\_fallback} silently switches that pair to point contact; and if mesh elements are much
larger than the contact surface, contact effectively becomes point contact.  Experiments therefore run the strict
\code{hydroelastic} model, which throws on a rigid--rigid pair (\code{drake\_config.py}, E0-2).  At
$E=10^{7}$ Pa the paddle contact penetrates about 2.5 $\mu$m, so a $0.002^\circ$ yaw misalignment puts the whole
patch at one edge: the paddle contact is modelled as a segment whose force centroid can move (Lynch and Mason), not
as an averaged face.  Drake combines a pair's friction as $2\mu_a\mu_b/(\mu_a+\mu_b)$; the object geometry carries a
fixed value and each strip and the paddle get the partner value that makes the pair coefficient equal to its target
(\code{scene.partner\_mu}), so one object can touch strips of different friction.

\section{Generative model}
\label{sec:generative}
\paragraph{Push map (release model).}  Contact holds while the paddle decelerates more slowly than friction would
decelerate the free object:
\begin{equation}
t^{*}=\min\{\,t>t_{\text{peak}}:\ a_p(t)<-\mu g\,\},\qquad v_r=v_p(t^{*}).
\label{eq:release}
\end{equation}
For braking that starts at zero deceleration and ramps at constant jerk $J$,
\begin{equation}
v_r=v_{\text{peak}}-\frac{(\mu g)^2}{2J}.
\label{eq:vr}
\end{equation}
At $J=100$ m/s$^3$ and $v_{\text{peak}}=0.8$ m/s this costs 11\% of the distance at $\mu=0.3$ and about 1\% at
$\mu=0.1$ (\code{physics.release\_speed}; \code{release\_from\_profile} evaluates~\eqref{eq:release} on any sampled
paddle profile).  Under A2--A3, heading $\psi_r\approx0$ and spin $\omega_r\approx0$ at release.

\paragraph{Free slide.}  With uniform friction,
\begin{equation}
x(\tau)=x_r+v_r\tau-\tfrac12\mu g\tau^{2},\quad 0\le\tau\le T=\frac{v_r}{\mu g},\qquad
x_{\text{rest}}=x_r+\frac{v_r^{2}}{2\mu g}.
\label{eq:slide}
\end{equation}
Across strips the deceleration is $\mu(x)g$, integrated piecewise until the speed reaches zero
(\code{physics.slide\_piecewise}, \code{slide\_field}).  Spin is a diagnostic only.  For a base with symmetric
pressure spin \emph{lengthens} the slide, because sliding and spinning share the Coulomb friction (Farkas et al.):
a uniform-pressure disk gives $d/d_0\approx1+0.26\,(R\omega_r/v_r)^2$ for $R\omega_r/v_r\le1$, with 1.065 at 0.5 and
1.27 at 1 from the integrated model (\code{physics.disk\_spin\_slide}).  A heading error $\psi$ shortens the
along-axis distance only by $\psi^2/2$ and adds a lateral miss of $\psi d$.

\paragraph{Transition across practice slides.}  The parameters do not change between slides:
\begin{equation}
\phi_{j+1}=\phi_j .
\end{equation}
Yu et al.\ measured plywood friction falling from 0.28 to 0.24 over 100 passes, about 0.15\% per pass, which is
negligible over five slides, and the object is reset between slides.  This is a decision, not an omission.

\emph{Added 17 Nov.}  The identity is about time.  It also assumes friction is the same along the board, since the
practice slide and the scored shot cross different stretches of it; that part was tested in Drake and does not hold on a
board whose friction varies from place to place (Section~\ref{sec:status-nov}, ``The transition model, tested'').

\section{Observation models}
\label{sec:observation}
\paragraph{$O_{\text{LLM}}$: the language model as a sensor.}  The prior $p_0(\theta)$ is the population
distribution of measured $\ln\mu$, fitted on training folds (\code{belief.prior\_from\_measured}).  The LLM's
median $m$ enters through a likelihood whose parameters are fitted leave-one-object-out on measured pairs
(\code{belief.fit\_llm\_likelihood\_loo}):
\begin{equation}
m_{\text{LLM}}=\alpha+\beta\theta+\varepsilon,\qquad\varepsilon\sim\mathcal N(0,s^{2}).
\end{equation}
This is the old ``recalibration'' step, now the observation model.  If $\beta\to0$ the LLM carries no
information.  The prompt builder reads only the \code{description\_*} and \code{photo\_*} columns of the pair file
(\code{llm\_prior.py}, test 7); prompts are frozen and responses cached before the real measured file exists.

\paragraph{$O_{\text{phys}}$, model C (primary).}  Each practice slide gives three measurements of~\eqref{eq:slide}:
frames $\hat x(t_i)$ of the object during the free slide from the depth camera and ICP, noise $\sigma_f$ per frame;
the rest pose, averaged over many stationary frames, noise $\sigma_{\text{rest}}$; and the release point, from arm
forward kinematics at separation, noise $\sigma_{\text{FK}}$:
\begin{equation}
\hat x_i=x(t_i;x_r,v_r,\mu)+e_i,\qquad
\hat x_{\text{rest}}=x_r+\frac{v_r^{2}}{2\mu g}+e_{\text{rest}},\qquad
\hat x_r^{\text{FK}}=x_r+e_{\text{FK}} .
\end{equation}
Model A uses frames only; model B adds the rest pose; model C adds the release point
(\code{observe.fit\_slide}).  Frames are timed from the arm's separation time, $t_i=i/f$ for $0\le t_i\le T$.

\paragraph{O1 (ablation only).}  Endpoint plus commanded speed, $y=\ln\big(v_{\text{cmd}}^2/2gd\big)$.  It carries the
release-ramp bias $(v_{\text{cmd}}/v_r)^2-1$, which grows with $\mu$ by~\eqref{eq:vr}.  It is kept only to show that
bias (\code{observe.fit\_o1}, \code{o1\_bias}).

\paragraph{Error budget.}  Relative standard error of $\mu$ from one slide (Cram\'er--Rao bound,
\code{observe.crlb}; test 2 reproduces every cell).  Camera at 60 Hz (decision 8, settled 5 Nov; the 31 Oct budget used 30 Hz) with 2 mm per-frame noise; rest pose and FK
at 0.5 mm.  O1 assumes 3\% speed error and 2 mm endpoint error; its bias is computed at braking jerk $J$.
\begin{center}\small
\begin{tabular}{@{}cclrrrrrr@{}}
\toprule
$\mu$ & $v_r$ (m/s) & Slide & A: frames & B: $+$rest & C: $+$FK & O1 noise & O1 bias, $J{=}100$ & O1 bias, $J{=}300$\\
\midrule
% generated by observe.py: relative standard error of mu from one slide (60 Hz, 2 mm frames; 0.5 mm rest and FK)
0.1 & 0.8 & 0.82 s, 32.6 cm & 1.2\% & 0.9\% & 0.7\% & 6.0\% & +1.2\% & +0.4\% \\
0.1 & 1.1 & 1.12 s, 61.7 cm & 0.5\% & 0.4\% & 0.3\% & 6.0\% & +0.9\% & +0.3\% \\
0.3 & 0.5 & 0.17 s, 4.2 cm & 16.7\% & 12.3\% & 9.0\% & 7.6\% & +19.9\% & +6.0\% \\
0.3 & 0.8 & 0.27 s, 10.9 cm & 5.6\% & 4.0\% & 2.9\% & 6.3\% & +11.8\% & +3.7\% \\
0.3 & 1.1 & 0.37 s, 20.6 cm & 2.6\% & 1.9\% & 1.4\% & 6.1\% & +8.4\% & +2.7\% \\
0.5 & 0.5 & 0.10 s, 2.5 cm & 32.0\% & 24.5\% & 18.6\% & 9.9\% & +73.4\% & +18.2\% \\
0.5 & 0.8 & 0.16 s, 6.5 cm & 12.8\% & 8.1\% & 6.0\% & 6.7\% & +38.5\% & +10.8\% \\
0.5 & 1.1 & 0.22 s, 12.3 cm & 5.4\% & 3.8\% & 2.8\% & 6.2\% & +26.1\% & +7.7\% \\

\bottomrule
\end{tabular}
\end{center}
Two consequences.  Model C is unbiased but its variance grows fast as slides get short: the error scales roughly
as $\sigma_f/(\sqrt N\,T^{2})$ with $T=v/\mu g$.  And the prior's worth in practice slides,
$n_0=\sigma_{\text{obs}}^{2}/s_0^{2}$ (\code{belief.n0}), now depends on the object: with $s_0=0.45$ it ranges from
about $10^{-4}$ slides for slippery objects to about 0.3 for $\mu=0.5$ at 0.5 m/s.

\section{Inference}
\label{sec:inference}
A grid of 500 points over $\theta\in[\ln0.05,\ln1.0]$ (\code{belief.theta\_grid}).  The physics likelihood is
Gaussian in $\mu$ while the prior and the LLM model are Gaussian in $\ln\mu$, so the grid is exact where a conjugate
update is not (test 4 checks it against the conjugate closed form for Gaussian-in-$\theta$ observations).  The
per-slide nuisance $(x_r,v_r)$ is profiled out by least squares at each grid point
(\code{belief.profile\_nuisance}, vectorized Gauss--Newton):
\begin{equation}
p(\theta\mid m,z_{1:k})\;\propto\;p_0(\theta)\,p(m\mid\theta)\prod_{j=1}^{k}\max_{x_r,v_r}p(z_j\mid\theta,x_r,v_r).
\end{equation}
At level L3 the grid becomes two-dimensional over $(\mu,e)$.  Pushing outcomes can be heteroscedastic and
multimodal (Bauza), which a single Gaussian cannot represent but a grid can.  The grid's floor at $\mu=0.05$ acts as a
hard prior: for a wide posterior near it, the $\mu^{-2}$ weight of the squared-relative loss makes the decision
depend on that floor (test 5 notes the effect).

\section{Decision rule}
\label{sec:decision}
Because $v_r$ depends on $\mu$, the closed-form release speed only holds when braking is abrupt.  The commanded peak
speed $u$ is chosen by a one-dimensional numerical search on the expected loss, with the release model inside $d$
(\code{decide.choose\_speed}):
\begin{equation}
u^{*}=\arg\min_{u}\int L\big(d(\mu,u)-\dstar\big)\,p(\mu\mid\text{data})\,d\mu,\qquad
d(\mu,u)=\frac{v_r(\mu,u)^{2}}{2\mu g}.
\end{equation}
The old closed form, $v_{\text{opt}}^{2}=2g\,\dstar\exp(m_k-1.5S_k^{2})$, is kept as the unit-test limit for
$J\to\infty$ with squared relative loss (test 5).  Losses: squared relative error, absolute relative error, a miss
outside $\max(0.1\dstar,\,1\text{ cm})$, and the table-edge penalty at $L=1.3\dstar$.  With $\ln\mu\sim\mathcal N(\ln0.3,0.3^2)$
and $\dstar=0.3$ m, the optimum moves from 1.242 m/s at $J=\infty$ to 1.253 at $J=300$ and 1.274 at $J=100$ m/s$^3$.

\textbf{Practice slides run at the maximum feasible speed} (\code{decide.practice\_speed}).  The standard error
falls roughly as $v^{-2.5}$, so the Fisher information grows roughly as $v^{5}$.  That is valid only if $\mu$ does not
depend on speed, which the rig's $\mu(v)$ data check (A4).  The slide room is the 48 in board's far edge less the release
point and the object's half length (\code{decide.slide\_room}; decision 1, 5 Nov; \code{scene.Table} tiles the same board,
$x = 0.3$ to $1.52$ m ahead of the iiwa base).  The release point is the start plus the push's travel to release
(\code{physics.release\_push\_distance}), so the room shrinks as the commanded speed grows.

\section{Fidelity ladder}
\label{sec:ladder}
\begin{center}\small
\begin{tabular}{@{}lp{3.0cm}p{2.9cm}p{2.6cm}p{2.2cm}p{4.0cm}@{}}
\toprule
Level & Unknowns & Observation & Inference & Experiments & Expected difficulty\\
\midrule
L0 & $\mu$ only, perfect measurements & $d$ and $v_r$ exact & closed form & pipeline check & trivial: one slide fixes $\mu$; the prior matters only at $k=0$\\
L1 & $+$ execution and perception noise & model C vs O1 & 1-D grid & E0-5, E2--E4 & O1's release-ramp bias; model C's variance at high $\mu$\\
L2 & $+$ friction field (strips) & model C per lane & grid with lane offsets & E0-6, E4 & which lane to practise on\\
L3 & $+$ COM offset $e$, point contact & $+$ landing yaw and lateral offset & 2-D grid $(\mu,e)$ & E0-3 & mostly lateral error; distance error second order\\
L4 & $+$ simulator vs analytic mismatch & as above & Drake rollouts as the likelihood & E0-1; Isaac branch & contact parameters, discrete approximation\\
\bottomrule
\end{tabular}
\end{center}
The residual of the constant-$\mu$ fit tells which level is active (\code{observe.residual\_signature}, by BIC over
closed-form alternatives): a clean parabola with residuals at noise level (constant $\mu$ on a uniform table); a
step in deceleration (a friction strip crossed); deceleration that changes smoothly (spin, as the spin-to-slide
ratio evolves; or a speed-dependent $\mu$, which across slides shows as $\hat\mu$ correlated with $v_r$,
\code{observe.speed\_signature}).  \emph{Detectability}: on synthetic 1.1 m/s slides with $\mu$ stepping from 0.1
to 0.2 part way, the step is labelled on 20 of 20 trials at 0.5 mm frame noise and 60 Hz and on 18 of 20 at 1 mm,
but at 2 mm and 30 Hz (the 31 Oct budget) it reads as a smooth trend on 20 of 20.  Telling a strip from spin in one slide
needs sub-millimetre frames or 60 Hz and up, which bears on the camera-rate decision (E0-5).

Prediction to report against: L0 is easy, as the instructor expects.  The real difficulty is the $\mu$-dependent
release and the simulator-vs-analytic mismatch, not friction or COM uncertainty as such.

\section{Corrections to earlier drafts}
\label{sec:corrections}
\begin{enumerate}\small
\item Off-centre friction does not make the path ``curve and shorten'': for a symmetric base, spin lengthens the slide
      ($d/d_0\approx1+0.26(R\omega/v)^2$); a heading error $\psi$ shortens the along-axis distance only by $\psi^2/2$ and
      adds a lateral miss $\psi d$.
\item A trajectory fit estimates $\mu$ to 1.7\% only for long slides; at $\mu=0.5$, $v=0.5$ m/s frames alone give 37\%.
\item The standard error falls as $v^{-2.5}$, so the information grows as $v^{5}$ (not $v^{2.5}$).
\item SAP ignores Hunt--Crossley dissipation and uses \code{relaxation\_time}; Lagged and Similar use Hunt--Crossley
      dissipation and ignore \code{relaxation\_time}.  Each parameter is paired with its approximation.
\item Release speed is not the paddle's peak speed: $v_r=v_{\text{peak}}-(\mu g)^2/2J$, and the error grows with $\mu$.
\item Lookup rows for Yu's pairs are not Yu's values: Yu sets simulator truth only (\code{ground\_truth\_published.csv});
      the lookup (\code{friction\_lookup.csv} v2) carries independent values for the same pair IDs.
\item The rig tests objects O1--O9 on surfaces S1--S3 as defined in protocol v3, not the materials listed in lookup
      table v1.
\end{enumerate}

\section{Implementation map and status}
\label{sec:impl}
\begin{center}\small
\begin{tabular}{@{}lp{7.6cm}p{5.2cm}@{}}
\toprule
Module & Functions & Tests (\code{pytest tests/})\\
\midrule
\code{physics.py} & \code{stop\_distance}, \code{release\_speed}, \code{release\_from\_profile}, \code{release\_push\_distance}, \code{slide\_trajectory}, \code{slide\_piecewise}, \code{slide\_field}, \code{spin\_factor}, \code{disk\_spin\_slide}, \code{no\_tip} & 1: sampled release vs~\eqref{eq:vr} to $10^{-6}$ m/s; 6: disk $d/d_0$ 1.065, 1.27\\
\code{observe.py} & \code{simulate\_obs}, \code{fit\_slide} (A, B, C), \code{fit\_o1}, \code{crlb}, \code{error\_budget}, \code{residual\_signature}, \code{speed\_signature} & 2: budget table to 0.1 pp; 3: 2,000 slides, bias $<1\%$, SD within 20\% of CRLB\\
\code{belief.py} & \code{theta\_grid}, \code{prior\_from\_measured}, \code{loglik\_lookup}, \code{lookup\_observation} (C2: the table's value as one observation of $\theta$, decision 3), \code{fit\_llm\_likelihood(\_loo)}, \code{loglik\_slide}, \code{posterior}, \code{n0} & 4: grid vs conjugate update to $10^{-3}$, with and without the lookup term\\
\code{decide.py} & \code{expected\_loss}, \code{choose\_speed}, losses, \code{closed\_form\_speed}, \code{practice\_speed}, \code{slide\_room} & 5: $J\to\infty$ limit to three significant digits\\
\code{llm\_prior.py} & \code{build\_prompt}, \code{cached\_response} & 7: only \code{description\_*}, \code{photo\_*} reach a prompt\\
\code{drake\_config.py}, \code{sim\_config.yaml} & \code{load\_config}, \code{contact\_from\_config}, \code{describe}, \code{assert\_hydro\_only} & E0-2 (\code{e0\_2\_hydro\_check.py})\\
\code{e0\_1\_free\_slide.py} & grid, spin and heading launches & E0-1 pass criteria; spin within 10\% of the disk model\\
\code{push.py} & \code{Pusher}, \code{profile}, \code{push\_slide}: a prescribed plate under Drake's implicit PD stands in for the arm; the release is read at $t^*$ on the command, on the recorded motion, and when the contact force vanishes & \code{tests/test\_push.py}: the command reproduces~\eqref{eq:vr}; one push within 1\%\\
\code{e0\_3\_contact\_geometry.py}, \code{e0\_4\_release.py}, \code{e0\_5\_observation.py} & E0-3, E0-4, E0-5 & their pass criteria; evidence for decisions 7 and 8\\
\bottomrule
\end{tabular}
\end{center}
\paragraph{E0 status (31 Oct).}  \emph{E0-2 passes}: in the passive and the push scene, for box, disc, ring and
footed objects, in strict and in fallback mode, every contact over 300 steps is hydroelastic (object--table every
step; paddle--object every step in the push scene) and there are zero point-pair contacts
(\code{e0\_2\_results.csv}, \code{e0\_2\_contacts.png}).  Negative controls: the paddle lowered 1 mm into the table
makes strict mode throw; a free rigid box on the rigid table runs silently on point contact in fallback mode.
\emph{E0-1 freezes A6} at the \code{sim\_config.yaml} values (Lagged, 1 ms, $E=10^{7}$ Pa, dissipation 1 s/m,
margin $10^{-4}$ m): box slides within 0.96\% of $d_0$ and $0.00^\circ$ of yaw over $\mu\in\{0.1,0.3,0.5\}$,
$v\in\{0.6,1.0\}$ m/s; disc spin launches within 2\% of the disk integrator (1.053--1.060 at $R\omega/v=0.5$,
1.249--1.266 at 1); heading launches within 1\% along and across with no yaw drift.  Finding: with the $10^{-4}$ m
margin, Hunt--Crossley dissipation of 10 s/m or more under Lagged (1 s/m or more under Similar) shortens slides by 4
to 85\%, while without margin Lagged passes up to 50 s/m; 2 ms steps pass only without margin, 5 ms never; SAP
passes at 1 ms with relaxation time 0.01 s or less (\code{e0\_1\_summary.txt}).  19 of 66 settings pass.

\paragraph{E0 status (5 Nov).}  The arm still has no controller, so \code{push.py} drives the paddle plate on a prismatic
joint under Drake's implicit PD actuator, tracking a profile that accelerates at 2 m/s$^2$ and then runs one jerk segment $-J$
through the peak into the brake, so that braking starts from zero deceleration at $v_{\mathrm{peak}}$ as \eqref{eq:vr} assumes.
The command is fed one time step ahead, after which the plate tracks to 0.05 mm.  Near separation both bodies decelerate at
about $\mu g$, so the instant chosen moves the speed by $\mu g\,\Delta t$ per step (1.4\% of a 0.36 m/s release at 1 ms); the
reported speed is the simulator's object state at the model's instant $t^* = t_{\mathrm{peak}} + \mu g/J$.
\emph{E0-3 (A2)}: with the plate covering the whole face, the box leaves with $|\psi_r| \le 0.004^\circ$ and a lateral miss
$\le 2$ mm for every COM offset up to $w/4$, plate yaw up to $2^\circ$ and $\mu_p \in \{0.1, 0.3, 0.6\}$; a yawed plate turns the box
to its own angle without a heading.  A disc with an offset COM spins off (heading 5--27$^\circ$, miss 23--116 mm at $e = R/4$, $R/2$)
and follows a yawed plate at about two thirds of its angle.  With the COM 5 mm inside the footprint edge the box rocks onto the edge
(tilt 4--5$^\circ$, heading to $6^\circ$) and every such disc topples; with the covered segment not straddling the COM the box
turns $90^\circ$ and leaves 33--86\% below the release model.  The slide error stays within 0.8\% wherever the object does not tilt.
\emph{E0-4 (A3)}: at the frozen contact the object's speed at $t^*$ is within 1\% of \eqref{eq:vr} in 53 of 54 cases over
$J \in \{30, \dots, 2000\}$ m/s$^3$, $\mu \in \{0.1, 0.3, 0.5\}$, $v_{\mathrm{peak}} \in \{0.6, 0.8, 1.0\}$ m/s (worst $+1.2\%$ at a
0.2 m/s release), and within 0.08\% of the model on the plate's recorded motion; the contact force vanishes 0--2 ms after $t^*$.
Under Similar and SAP, with any damping, the box tilts 0.4--1.2$^\circ$ during the push and leaves 1--7\% off the model with spin:
the release model needs Lagged, which supports A6.  The distance lost to the ramp (decision 7) is tabulated in
\code{e0\_4\_summary.txt}: under 2\% needs $J \approx 500$ m/s$^3$ at $\mu = 0.3$ and about 2000 at $\mu = 0.5$.
\emph{E0-5}: pushes replayed through the rendered depth camera at 30, 60 and 120 Hz, 20 replicates each, give models A, B and C
within 1\% of the simulator's $\mu$ with standard errors at or below the Cram\'er--Rao bound at the measured noise (ratio
$\le 0.73$); the residual $-0.6$ to $-0.9\%$ bias at $\mu = 0.3$--$0.5$ is the first frame after release, taken against the plate
(dropping the first 30 ms leaves $\le 0.3\%$).  O1's bias matches $(v_{\mathrm{peak}}/v_r)^2 - 1$.  Measured noise for Tier B:
$\sigma_f = 0.09$ mm, $\sigma_{\mathrm{rest}} = 0.03$ mm near the camera axis and 0.1 mm at the far end of the table; $\sigma_{\mathrm{FK}}$ is
assumed 0.5 mm because the simulator's forward kinematics are exact.  The Section~\ref{sec:observation} budget assumed 2 mm frames and
is therefore 10--20$\times$ pessimistic; the measured values are lower bounds (exact camera pose and intrinsics, no motion blur).
Decision 8 is not settled by $\mu$'s precision (model C: 0.42\% at 30 Hz and 0.25\% at 120 Hz on the shortest slide) but by what
else the frames must resolve.  The team chose 60 Hz on 5 Nov (\code{observe.CAMERA\_FPS}); the
error-budget table of Section~\ref{sec:observation} is at that rate.

\section{Status, 6--17 Nov: the arm, the comparison, and what the results say about the model}
\label{sec:status-nov}
Sections~\ref{sec:state}--\ref{sec:ladder} were written before the arm pushed and before the method comparison ran.
This section records what changed and what was found; the README and \code{results/side\_base/README.md} have the
tables.  All numbers are from the simulator.

\paragraph{The arm (6 and 16 Nov).}  \code{arm.ArmPusher} replaces the prescribed plate.  The push is a straight line, so
the paddle pose fixes link 7's pose and leaves one redundant joint; the start posture is the multi-start IK solution whose
path, integrated with the minimum-norm joint tangent, has the smallest peak joint-speed ratio inside the joint limits.
Each joint runs Drake's implicit PD on the desired trajectory with the arm's inverse dynamics as feed-forward torque.  The
brake ramps at jerk $J$ to a cap (12 m/s$^2$ by default) and holds it; the release model is unchanged while the cap
exceeds $\mu g$.  At $\mu=0.3$, 0.8 m/s, $J=100$ the object leaves at $+0.11\%$ of~\eqref{eq:vr} and slides to within
$-0.20\%$ of $v_r^2/2\mu g$.  Over $\mu\in[0.1,0.5]$ and 0.3--1.0 m/s the release speed is within $-0.02$ to $+1.23\%$ of
the model (\code{experiments/e1\_release\_error.py}).  With the base on the push line joint 4 limits the push to 0.49 m/s;
with the base beside the board at $(0.6, 0.65)$ m (the default since 16 Nov, \code{scene.LAYOUTS}) joint 1 sweeps the
paddle and the bound is 1.02 m/s inside the joint speed and torque limits.  The joint acceleration limits are met only
with a softer brake (4 m/s$^2$, which releases $\mu<0.4$).

\paragraph{Point targets (16 Nov).}  The decision rule of Section~\ref{sec:decision} is stated for a slide $\dstar$ from the
release point.  A target \emph{point} $x^\star$ on the board is the same rule with a slide that moves with the speed,
\begin{equation}
\dstar(u,\mu)=x^\star-x_r(u,\mu),\qquad x_r=x_0-\text{gap}+\text{(push travel to release)},
\end{equation}
because a faster push carries the object further before it lets go (\code{decide.point\_target},
\code{arm.plan\_peak\_speed}).  The speed bound and the targets are computed from the layout, not hand-set: the bound is the
fastest push inside the checked joint limits, and the targets run from $x_0+5$ cm to what the highest-friction pair reaches
at 90\% of the bound (\code{arm.limited\_peak\_speed}, \code{arm.target\_grid}).

\paragraph{The simplest version (the instructor's starting point).}  Friction the only unknown, no measurement noise,
pairs Y01--Y03, three targets.  Before any practice the generic, lookup and LLM priors miss by 39.8, 41.3 and 48.4 mm
on average.  After one practice slide all three miss by 0.15--0.16 mm, the same as a plan made with the true friction
(0.26 mm).  It is straightforward, as predicted.

\paragraph{How much a practice slide says (17 Nov).}  One slide at the practice speed (1.01 m/s) gives $\ln\mu$ to 0.006
on its own at the nominal noise of Section~\ref{sec:observation}; the priors' standard deviations are 0.80 (generic),
0.33--0.51 (lookup) and 0.26--0.30 (LLM).  In slides, the LLM prior is worth about $(0.006/0.26)^2\approx 1/2000$ of
one.  \code{evaluation/practice\_study.py} weakens the slide (practice speed 0.15--1.01 m/s, noise 0.5--16 times
nominal): the methods' shots land more than 1 mm apart only once a slide alone leaves $\ln\mu$ uncertain by about 0.035,
which needs a practice push of 0.45 m/s or slower, or 8 times the nominal noise.  The camera as measured in E0-5 is about
12 times better than nominal.  The comparison between priors therefore lives in the shot made before any practice.

\paragraph{The board's seams (17 Nov).}  \code{push\_slide}'s uniform table is 24 equal-friction strips.  As the object
crosses a seam the contact patch vanishes for one time step and the deceleration spikes; about a third of pushes on
$\mu=0.14$ land short, by 1.3 mm RMS and at most 7.5 mm, and which pushes are hit changes with the speed to 0.01 m/s.  On a
one-piece board the landing is within 0.5 mm of the release model.  This is a simulator artefact, not friction; it is all
of the largest landing errors recorded so far.  Not changed, pending a team decision.

\paragraph{More pairs and more than one LLM (17 Nov).}  Until the rig data exists, the 20 handbook pairs of the lookup
table (metals, engineering plastics, oak, glass, graphite; $\mu$ 0.14--0.49) are a second truth set
(\code{data/ground\_truth\_handbook.csv}); the lookup baseline holds the answer on them, and an LLM may recall the same
handbooks.  First shot, mean miss: generic 61.3 mm, lookup 3.9 mm, LLM 13.7 mm (a Claude Haiku stand-in; GPT is cached for
Y01--Y03 only).  Most of that gain is not knowledge of the pair.  The generic prior's median is about as close to the truth
as the LLM's (RMS error of $\ln\mu$ 0.37 against 0.33); it is its width, 0.08--1.1, that makes the planner push short.  A
generic prior with the LLM's own interval width, pair by pair, misses by 19.7 mm (\code{evaluation/compare\_llms.py}), so
the steps are 61 mm (wide generic), 20 mm (equally confident generic), 14 mm (LLM).  On Y01--Y03, which no table lists,
neither LLM beats the generic prior.  After one practice slide all methods agree to 0.15 mm on all 23 pairs.

\paragraph{The transition model, tested (17 Nov).}  \code{evaluation/patchy\_table\_study.py} runs the comparison on the
scene's patchy board: 24 strips of 2 in with friction $\mu_i=\mu\,e^{s z_i}$, $z_i\sim\mathcal N(0,1)$, one draw per board,
the same board for the practice slide and the scored shots; pairs Y01--Y03, 16 boards per pair and spread, nominal noise.
Let $\mu_{\text{shot}}$ be the uniform friction under which the scored push would have landed where it did
(\eqref{eq:vr} and the free slide, inverted) and $\hat\mu$ the estimate from the practice slide.  The identity transition
says $\ln(\mu_{\text{shot}}/\hat\mu)=0$ to within the posterior's width, about 0.009.
\begin{center}\small
\begin{tabular}{@{}lcccc@{}}
\toprule
strip spread $s$ & 0 & 0.06 & 0.12 & 0.25\\
\midrule
SD of $\ln(\mu_{\text{shot}}/\hat\mu)$, practice at 1.01 m/s & 0.008 & 0.047 & 0.093 & 0.197\\
shots inside the belief's 90\% interval & 100\% & 19\% & 12\% & 8\%\\
mean miss, one practice slide at 1.01 m/s (mm) & 0.8 & 4.0 & 8.0 & 16.5\\
mean miss, one practice slide at 0.7 m/s (mm) & 0.6 & 3.5 & 5.9 & 11.8\\
mean miss, plan with the nominal $\mu$ (mm) & 0.3 & 2.7 & 5.5 & 11.9\\
mean miss, no practice, generic prior (mm) & 39.8 & 40.6 & 41.5 & 44.2\\
\bottomrule
\end{tabular}
\end{center}
The practice estimate and the friction the shot meets differ by about $0.8\,s$, so the belief after practice is
overconfident by an order of magnitude at the spreads of Yu et al.\ (their spatial SD over the mean: about 0.09 for plywood, 0.11 for acetal and 0.12 for ABS; 0.22 for polyurethane, which is not simulated.  \emph{Corrected 25 Nov:} this read ``0.11--0.25'', from a 0.25 in \code{ground\_truth\_published.csv} that was polyurethane's SD entered on the plywood row (fixed to 0.09 the same day); $s=0.25$ is about twice the largest spread of the simulated surfaces).  Practice still removes most of the
miss, but the floor is the board, not the sensor.  Where the practice slide runs matters: at 1.01 m/s it crosses
$x=0.83$--$1.19$ m while the shots cross $0.57$--$0.94$ m, and the result is worse than a plan made with the pair's nominal
friction; a practice slide at 0.7 m/s, on the stretch the shots use, matches the nominal plan.  This qualifies the rule of
Section~\ref{sec:decision} that practice runs at the maximum feasible speed, which was derived for uniform friction.
The three methods still make the same shot.

A transition with spread, $\theta_{\text{slide}}=\theta_{\text{pair}}+\eta$, $\eta\sim\mathcal N(0,\sigma^2)$ for the
practice slide and again for the shot (\code{practice\_study.update\_with\_transfer}; $\sigma$ set so that the
practice-to-shot spread is the measured one), makes the belief honest: 88--92\% of shots fall inside its 90\% interval.
It does not improve the shot (7.8--7.9 mm against 8.0 mm at $s=0.12$), because a wider belief cannot say which way the
shot's stretch differs from the practice stretch.  Told too large a spread it lets the priors back in and the miss grows.
The identity transition is therefore kept for planning; what it gets wrong is the stated uncertainty, and what would reduce
the miss is practising on the stretch the shot will cross.

\paragraph{Practice aimed at the target (17 Nov).}  That last statement is tested by
\code{patchy\_table\_study.py --aimed} on the same boards: the practice push is the method's own first shot at the target,
observed as a practice slide; the next shot at that target is planned from it, and a third from both slides (identity
transition) or from the second alone.
\begin{center}\small
\begin{tabular}{@{}lcccc@{}}
\toprule
strip spread $s$ & 0 & 0.06 & 0.12 & 0.25\\
\midrule
one aimed slide, generic / lookup / LLM (mm) & 3.4 / 4.2 / 2.8 & 3.9 / 3.8 / 3.8 & 5.1 / 5.6 / 6.0 & 9.4 / 10.0 / 11.0\\
two aimed slides (mm) & 2.1 / 2.6 / 1.9 & 2.4 / 2.6 / 3.0 & 3.0 / 3.9 / 4.5 & 5.5 / 7.3 / 8.4\\
two aimed slides, latest only (mm) & 2.7 / 2.5 / 2.1 & 3.0 / 3.1 / 2.9 & 4.0 / 3.7 / 3.9 & 7.3 / 6.2 / 6.4\\
fixed practice at 0.7 m/s (mm) & 0.6 & 3.5 & 5.9 & 11.8\\
plan with the nominal $\mu$ (mm) & 0.3 & 2.7 & 5.5 & 11.9\\
shots inside the 90\% interval, one aimed slide & 78\% & 83\% & 78\% & 58\%\\
\bottomrule
\end{tabular}
\end{center}
One aimed slide equals the slower fixed slide ($-0.2\pm0.4$ mm at $s=0.12$, paired by board).  Two aimed slides are the
first rule to beat the plan made with the pair's nominal friction, by $1.7\pm0.5$ mm at $s=0.12$ and $4.8\pm0.9$ mm at
$s=0.25$: the second slide ends within millimetres of the target and crosses the strips the shot will cross.  At the 0.77
and 0.94 m targets the SD of $\ln(\mu_{\text{shot}}/\hat\mu)$ at $s=0.12$ falls from 0.099 and 0.076 (practice at 1.01 m/s)
to 0.029 and 0.025, about $s/4$ in place of $0.8\,s$.  The price is information: a slide aimed at the 0.60 m target is a
few centimetres long and leaves a posterior SD of about 0.27 in $\ln\mu$, and on the uniform board the aimed rule misses
by 2--4 mm where a fixed practice slide gives 0.8 mm.  So Section~\ref{sec:decision}'s rule (practise fast) and this one
(practise where the shot goes) trade information against relevance, and which wins depends on $s$.

With the less informative slide the prior matters after practice: the three methods' shots land 5--19 mm apart.  With one
aimed slide at $s=0.12$ the miss is 3.4 / 6.6 / 7.5 mm on Y01 ($\mu=0.14$; lookup and LLM medians 0.26 and 0.28 with
narrower intervals) and 8.2 / 4.1 / 3.2 mm on Y03 ($\mu=0.26$; medians 0.25).  Over the three pairs, with two slides under
the identity transition, the LLM method is worse than the generic prior by $1.5\pm0.4$ mm ($s=0.12$) and $2.8\pm0.7$ mm
($s=0.25$).  Planning from the latest slide alone removes the difference ($-0.1\pm0.3$ mm), since the first slide, taken
far from the target, is the one that carries the prior's error.  Not compared: two slides at a fixed speed.

\paragraph{The aimed rule on 23 pairs (17 Nov).}  \code{patchy\_table\_study.py --aimed-pairs} repeats this on the 20
handbook pairs and Y01--Y03 (8 boards per pair and spread, targets 0.60, 0.735 and 0.87 m, 22{,}287 Drake shots), with the
Claude Haiku stand-in as the LLM and a control: the generic median with the LLM's interval width, pair by pair.
\begin{center}\small
\begin{tabular}{@{}lccc@{}}
\toprule
strip spread $s$ & 0 & 0.12 & 0.25\\
\midrule
no practice, generic / lookup / LLM / control (mm) & 57.7 / 7.9 / 17.3 / 22.6 & 57.7 / 8.6 / 17.5 / 22.6 & 58.0 / 11.0 / 18.4 / 23.2\\
one aimed slide (mm) & 6.3 / 2.6 / 2.2 / 2.9 & 7.3 / 3.0 / 3.8 / 4.1 & 10.1 / 4.4 / 5.5 / 6.1\\
two aimed slides (mm) & 4.6 / 1.8 / 1.6 / 2.4 & 4.4 / 2.4 / 3.0 / 3.1 & 6.0 / 3.8 / 4.6 / 4.8\\
fixed practice at 1.01 m/s, generic (mm) & 0.9 & 5.2 & 10.1\\
plan with the nominal $\mu$ (mm) & 0.2 & 3.6 & 7.3\\
\bottomrule
\end{tabular}
\end{center}
Paired over pairs, after one aimed slide at $s=0.12$ the lookup is better than the generic prior by $4.3\pm0.8$ mm and the
LLM by $3.6\pm0.8$ mm, but the control is better by $3.2\pm0.7$ mm: the LLM adds $0.4\pm0.3$ mm to the control and the
lookup, which on the handbook pairs holds the simulator's friction, $1.1\pm0.3$ mm.  Most of the gain is the prior's
width.  The set's frictions (0.14--0.49) straddle the generic median, which is almost as near a pair's friction as the
LLM's answer (mean $|\ln(\text{median}/\mu)|$ 0.34 against 0.29).  The wide prior loses at the 0.60 m target, where the
aimed slide leaves it a posterior SD of 0.50 in $\ln\mu$ (0.20 for the others) and a miss of 8.8 mm against 1.6--2.5 mm; at
the farther targets after two slides the methods are level.  With a confident prior, one aimed slide beats fixed-speed
practice on a patchy board and is level with the nominal-friction plan; on the uniform board fixed practice stays better.

\appendix
\section{sim\_config.yaml}
\begin{verbatim}
# sim_config.yaml -- the one Drake contact configuration that every experiment loads (drake_config.py).
# Status: FROZEN by E0-1 on 31 Oct 2026 (e0_1_free_slide.py; e0_1_summary.txt, e0_1_summary.png).  Change only with a new E0-1 run.
#   E0-1 result at these settings: box free slide within 0.96% of d = v^2/(2 mu g) and 0.00 deg of yaw over
#   mu in {0.1, 0.3, 0.5}, v in {0.6, 1.0} m/s; disc spin launches within 2% of the disk model; heading launches pass.
#   Interaction to respect: with margin 1e-4 m, Hunt-Crossley dissipation >= 10 s/m (lagged) or >= 1 s/m (similar)
#   shortens slides by 4-85%; without margin, lagged passes up to 50 s/m.  Keep dissipation <= 1 s/m with this margin.
#   2 ms steps pass only without margin (lagged) and 5 ms steps never pass.  SAP passes at 1 ms with relaxation <= 0.01 s.
time_step: 0.001                         # s                          (A6)
discrete_contact_approximation: lagged   # lagged | similar | sap     (A6, tentative)
contact_model: hydroelastic              # strict: a rigid-rigid pair throws instead of silently falling back to point contact
hydroelastic_modulus: 1.0e7              # Pa, object (compliant); Drake's default; swept 1e5-1e7 in e0_sweep.py
dissipation: 1.0                         # s/m, Hunt-Crossley: used by lagged and similar, IGNORED by sap
relaxation_time: null                    # s, sap only (null = Drake's default); IGNORED by lagged and similar
margin: 1.0e-4                           # m, hydroelastic margin on every geometry
resolution_hint: 0.01                    # m, object mesh resolution
compliance:                              # one compliant geometry per contact pair
  object: compliant
  table: rigid
  paddle: rigid
paddle_mu: 0.3                           # paddle-object pair coefficient; swept 0.1-0.6 in E0-3
friction_rule: >
  Drake combines a pair as 2 mu_a mu_b / (mu_a + mu_b).  The object geometry carries OBJECT_MU and
  scene.partner_mu gives each strip and the paddle the value that makes the PAIR coefficient equal to
  its target, so one object can touch strips of different friction.
\end{verbatim}

\end{document}
